\documentclass[10pt,a4paper]{amsart}
\usepackage[margin=19mm]{geometry}
\usepackage{amsmath,amssymb,mathtools}
\usepackage[scr=rsfs,cal=euler]{mathalfa}
\usepackage{xfrac}
\usepackage[unicode,hidelinks,bookmarks=false]{hyperref}
\newcommand{\ie}{i.e.\ }
\newcommand{\cf}{cf.\ }
\newcommand{\resp}{respectively\ }
\newcommand{\Subsection}{\subsection}
\providecommand{\nobreakdash}{\nobreak\hbox{-}}
\setlength{\emergencystretch}{2em}
\multlinegap0pt
\usepackage{bm}
\usepackage{bbm}
\usepackage{dsfont}
\usepackage{xspace}
\usepackage{cancel}
\usepackage{mathtools}
\usepackage{amsthm}
\makeatletter
\@ifundefined{proposition}{\newtheorem{proposition}{Proposition}}{}
\makeatother
\DeclareMathOperator{\re}{Re}
\title[An extension of smooth numbers: multiple dense divisibility]{An extension of smooth numbers: multiple dense divisibility}
\author{Garo Sarajian, Andreas Weingartner}
\date{Source: arXiv:2504.02699v2 (CC BY 4.0)}
\begin{document}
\maketitle

\noindent\emph{Source and licence.} An extension of smooth numbers: multiple dense divisibility. Version arXiv:2504.02699v2 (CC BY 4.0), distributed under the Creative Commons Attribution 4.0 International licence, as recorded in the arXiv licence metadata for this exact version and shown on the abstract page https://arxiv.org/abs/2504.02699v2. Full rights notice: licenses/arxiv-2504.02699-CC-BY-4.0.txt.\par\smallskip

\noindent\textit{Editorial scope.} Counted unit: the Proposition whose source label is \texttt{proplambdaub} in the article (source0.tex lines 1621--1628, source label \texttt{proplambdaub}), delivered together with the article's own complete original proof of it (source0.tex lines 1653--1813, one \texttt{proof} environment of 161 lines). No mathematics is written, repaired or re-derived: statement and proof are the article's own text line for line. The proof is detached from the statement in the source: the article states the result at lines 1621--1628 and prints its proof at lines 1653--1813, and the proof environment's own header names the statement, so the association is the article's own. Required context retained uncounted: gacint (lines 1401--1404). Every definition, display and label of those lines is retained unchanged. Omitted from this package and not redistributed: 1--1400, 1405--1620, 1629--1652, 1814--1898. Nothing inside a retained excerpt is omitted or altered: every retained body is delivered exactly as the article's lines are written, so no omission receipt of any kind accompanies this package. Citations: the retained excerpts carry no citation command at all, so no bibliography entry of the article is transcribed and no reference is redistributed. Reference adaptation: none was needed and none was made; every reference inside the delivered unit resolves inside it, so no unresolved reference or citation is produced. Printed slips kept literal: no printed word, symbol, hypothesis or number of the counted statement or of its proof was altered. Layout and class: the article's own class is \texttt{amsart} with packages including amsmath, amsfonts, amssymb, graphicx, enumerate, url; the wrapper is the lane's pinned \texttt{amsart} editorial wrapper at 10pt on A4, which restates the theorem-like environments the retained text needs and adds the article's own macro definitions that the retained text needs (\texttt{\textbackslash re}), with extra wrapper packages (bm, bbm, dsfont, xspace, cancel, mathtools). The delivered numbering is the unit's own. Assets: no retained span contains a figure environment, a graphics-inclusion command, a PDF-page inclusion, a table or a TikZ picture; no image file is copied into this package, the package declares no asset dependency and no image file is redistributed with it. Licence: version arXiv:2504.02699v2 of this article is distributed under the Creative Commons Attribution 4.0 International licence, as recorded in the arXiv licence metadata for this exact version and shown on the abstract page https://arxiv.org/abs/2504.02699v2. A full rights notice is delivered at licenses/arxiv-2504.02699-CC-BY-4.0.txt. Characters: every retained excerpt is pure ASCII, so the delivered engine, XeLaTeX with its default Latin Modern text font, reports no missing character.
\begin{equation}\label{gacint}
 g_a(-n) a e^\gamma 
= \frac{n!}{2\pi } \frac{e^{r-I(ar)}}{r^n}  \int_{-\pi}^\pi \exp\left(M_{a,r,n}(t)\right)dt,
\end{equation}

\begin{proposition}\label{proplambdaub}
Let $\nu>0$ be fixed. With the notation as above, we have 
\begin{equation*}
 g_a(-n) a e^\gamma 
= \frac{n!}{2\pi } \frac{e^{r-I(ar)}}{r^n}  \frac{\sqrt{2\alpha}}{L} ( K(\nu) + o(1)),
\end{equation*}
as $a \to 0$. 
\end{proposition}

\begin{proof}[Proof of Proposition \ref{proplambdaub}]
We approximate the exponent of the integrand in \eqref{gacint} by its fourth degree Taylor polynomial around $t=0$.
Here the linear term does not vanish.
We have
\begin{multline*}
M_{a,r,n}(t) = i  \left(r+1-n-e^{ar}\right)t
+\frac{r}{2}  \left(a e^{a r}-1\right) t^2
+ \frac{i r}{6} \left(a  e^{a r} (a r+1)-1\right) t^3 \\
+ \frac{r}{24}\left(1-a  e^{a r} \left(a^2 r^2+3 a r+1\right)\right) t^4 
+O(e^{ar} (ar)^4 |t|^5) \\
=: ic_1 t+ c_{2} t^2+ i  c_{3} t^3+  c_{4} t^4 +O( a^{-1} L^4 |t|^5),
\end{multline*}
since $\frac{d^5}{dt^5} M_{a,r,n}(t) \ll e^{ar} (ar)^4 \ll a^{-1} L^4$ for all $t\in \mathbb{R}$.
The coefficients are
$$
c_1 = r+1-n-e^{ar}= -\frac{\alpha}{a} +O(1),
$$
$$
c_2= \frac{r}{2}  \left(a e^{a r}-1\right)
=\frac{\alpha L}{6a}\left(1+O\left(\frac{1}{L}\right)\right),
$$
$$
c_3= \frac{r}{6} \left(a  e^{a r} (a r+1)-1\right) = \frac{L^2}{6a} ( 1+ O(\alpha)),
$$
and
$$
c_4 =  \frac{r}{24}\left(1-a  e^{a r} \left(a^2 r^2+3 a r+1\right)\right) = -\frac{L^3}{24 a} \left(1+O\left(\frac{1}{L}\right)\right).
$$

Define
$$
T=\frac{\alpha^{1/3}}{L^2}, \qquad T_1 = \frac{\alpha^{2/5}}{L^2}.
$$
For $|t|\le T$, we have $a^{-1}  L^4 |t|^5 \ll 1$. Thus,
\begin{equation*}
\begin{split}
 I_1 & := \int_{-T}^T \exp\{ M_{a,r,n}(t) \} dt   \\
& =\int_{-T}^T \exp\{c_2 t^2 + c_4 t^4\} \exp\{i(c_1 t + c_3 t^3)\} (1+O(L^4 a^{-1} T^5)) dt \\
& = O(L^4  a^{-1}  T^6)+ \int_{-T}^T \exp\{c_2 t^2 + c_4 t^4\}\exp\{i(c_1 t + c_3 t^3)\} dt, \\
& = O(\sqrt{\alpha}/L^7)+ \int_{-T}^T \exp\{c_2 t^2 + c_4 t^4\} \cos(c_1 t + c_3 t^3) dt, \\
\end{split}
\end{equation*}
since $ \exp\{c_2 t^2 + c_4 t^4\} \ll 1$ for all $t\in \mathbb{R}$.
Indeed, $c_2 t^2 + c_4 t^4$ has its global maximum at $t=\pm T_0$, say, where $T_0 \sim  \sqrt{2\alpha}/L$ and 
$c_2 T_0^2 +c_4 T_0^4 \ll 1$. 
Since $ \exp\{c_2 t^2 + c_4 t^4\}$ is decreasing on $[T_0, \infty]$ and $T_1>T_0$, the second mean value theorem for integrals 
implies that there exists a $T_2 \in [T_1,T]$ such that
$$
\int_{T_1}^T \exp\{c_2 t^2 + c_4 t^4\} \cos(c_1 t + c_3 t^3) dt
= \exp\{c_2 T_1^2 + c_4 T_1^4\} \int_{T_1}^{T_2} \cos(c_1 t + c_3 t^3) dt.
$$
A little exercise shows that the last integral is $\ll \frac{1}{c_3 T_1^2}$, so that 
$$
\int_{T_1}^T \exp\{c_2 t^2 + c_4 t^4\} \cos(c_1 t + c_3 t^3) dt \ll \frac{ 1}{c_3 T_1^2}\asymp \alpha^{7/10}L .
$$
When $|t| \le T_1$, we have
$$
|c_2 t^2 + c_4 t^4| \le |c_2 T_1^2 + c_4 T_1^4|  \asymp |c_4 T_1^4| \ll 1,
$$
which yields
$$
\exp\{c_2 t^2 + c_4 t^4\} = 1+ O( |c_4| T_1^4) 
$$
and
$$
 \int_{-T_1}^{T_1} \exp\{c_2 t^2 + c_4 t^4\} \cos(c_1 t + c_3 t^3) dt = O(|c_4| T_1^5)+ \int_{-T_1}^{T_1} \cos(c_1 t + c_3 t^3) dt .
$$
Note that $|c_4| T_1^5 \asymp \sqrt{\alpha} L^{-6}$. 
After the change of variable $u=t L/\sqrt{2 \alpha}$ and $U=T_1 L/\sqrt{2 \alpha} \asymp \alpha^{-1/10} /L$, 
 we have
\begin{equation*}
\begin{split}
&\frac{L}{\sqrt{2\alpha}}\int_{-T_1}^{T_1} \cos(c_1 t + c_3 t^3) dt\\
& =  \int_{-U}^U \left[\cos\left( -\frac{3\nu}{2}u +\frac{\nu}{2}u^3 \right) +O(U\sqrt{\alpha}/L+U^3 \alpha)\right] du\\
 & = O(  U^2 \sqrt{\alpha}/L + U^4 \alpha)+  \int_{-U}^U \cos\left( -\frac{3\nu}{2}u +\frac{\nu}{2}u^3 \right) du\\
 & = O(U^2 \sqrt{\alpha}/L  + U^4 \alpha +U^{-2})+  \int_{-\infty}^\infty \cos\left( -\frac{3\nu}{2}u +\frac{\nu}{2}u^3 \right) du\\
 &  = o(1)+  K(\nu).
\end{split}
\end{equation*}
Combining everything, we obtain
$$
  \int_{-T}^T \exp\{ M_{a,r,n}(t) \} dt = \frac{\sqrt{2\alpha}}{L} ( K(\nu) + o(1)),
$$
as $a \to 0$. 

It remains to show that 
$$
 \int_{T}^{\pi} \exp\{ M_{a,r,n}(t) \} dt = o\left(\frac{\sqrt{2\alpha}}{L} \right).
$$
Let $t_0 = 7/L$ and write
$$
 \int_{T}^{\pi} \exp\{ M_{a,r,n}(t) \} dt = \int_{T}^{t_0} + \int_{t_0}^{\pi} =: I_2 +I_3.
$$
For $I_2$ and $I_3$, we estimate the absolute value of the integrand. As in the case of a real saddle point,  we have
$$
\re(M_{a,r,n}(t)) = -r(1-\cos t ) + Q_1+Q_2,
$$
where
\begin{equation*}
Q_1+Q_2 =  \int_0^1 e^{arv}(1-\cos(arv \sin t) ) \frac{dv}{v} +  \int_0^1 ((e^{arv}-e^{arv \cos t })\cos(arv \sin t) ) \frac{dv}{v}.
\end{equation*}
We note that 
$$\frac{t^2}{5} \le 1-\cos t \le  \frac{t^2}{2}\qquad (|t| \le \pi).
$$ 
To estimate $I_3$, we write
\begin{multline*}
Q_1 \le \frac{1}{2}\int_0^{1/(ar t)} e^{arv}( a r  t)^2 v dv +  \int_{1/(ar t)}^{1/2} e^{arv} 2 \frac{dv}{v} + \int_{1/2}^1 e^{arv} 2 \frac{dv}{v}\\
\le \frac{t}{2}e^{1/t} + 2 (ar t) \frac{e^{ar/2}}{ar} + 4 \frac{e^{ar}}{ar}  \le  \frac{5r}{L^2},
\end{multline*}
for $a$ sufficiently small and $|t|\ge t_0$.
Since $1-e^{-x} \le x$ for all $x$, 
\begin{multline*}
Q_2 \le \int_0^{1} e^{arv}( 1-e^{-arv t^2/2}) \frac{dv}{v}\le \int_0^{1/2} e^{arv} ar (t^2/2) dv +  \int_{1/2}^1  e^{arv} 2 dv \\
\le e^{ar/2}\pi^2/2 + \frac{2 e^{ar}}{ar} \le \frac{3r}{L^2}.
\end{multline*}
We have $r(1-\cos t) \ge r t^2/5 \ge r t_0^2/5$. Combining these three estimates we obtain
$$
|I_3| \le 2\pi \exp( - rt_0^2/5 +8r/L^2) \le 2\pi \exp(-r/L^2)= o\left(\frac{\sqrt{2\alpha}}{L} \right).
$$

Next, we estimate $I_2$. For all real $t$ we have
$$
 r(1-\cos t) \ge r \left(\frac{ t^2}{2}-\frac{t^4}{24}\right).
$$

When $|t| \le t_0 = 7/L$, then $|\sin t| \le 7/L$ and 
$$
1-\cos(arv \sin t) \le \frac{(arv t)^2}{2}- \frac{(arvt)^4}{\kappa},
$$
for some suitable constant $\kappa$. Thus 
$$
Q_1 \le \int_0^1 e^{arv} \left(\frac{(arv  t)^2}{2} -\frac{(arvt)^4}{\kappa}\right) \frac{dv}{v}
\le  \frac{t^2}{2} \left( ar e^{ar} - e^{ar}+1\right)- \frac{t^4}{2\kappa} (ar)^3 e^{ar}
$$
and
$$
Q_2 \le \int_0^{1} e^{arv}( 1-e^{-arv t^2/2})  \frac{dv}{v} \le  \int_0^{1} e^{arv} \frac{ar t^2}{2} dv = \frac{t^2}{2} (e^{ar}-1).
$$
For sufficiently small values of $a$ and  $t_0\ge |t| \ge T = \frac{\alpha^{1/3}}{L^2}$, we obtain
\begin{equation*}
\begin{split}
  -r(1-\cos t) + Q_1+Q_2 
 \le & -r\left(\frac{ t^2}{2}-\frac{t^4}{24}\right)+ \frac{t^2}{2} ar e^{ar}  - \frac{t^4}{2\kappa} (ar)^3 e^{ar}\\
\le & \frac{t^2}{2} r \left(-1 + a e^{ar} \right) - \frac{t^4}{3\kappa} (ar)^3 e^{ar}\\
 \le & \frac{t^2}{2} r \alpha - \frac{t^4}{4\kappa} r L^2 
\le  - \frac{t^4}{5\kappa} r   L^2 
\le  -\frac{T^4}{5\kappa} r   L^2 \\
= &  -\frac{\alpha^{4/3}}{5\kappa L^6} r \le - a^{-1/10}.
\end{split}
\end{equation*}
It follows that 
$
|I_2| \le 2 \pi \exp(- a^{-1/10}) \ll a  = o\left(\frac{\sqrt{2\alpha}}{L} \right).
$
Combining everything, we have 
$$
  \int_{-\pi}^\pi \exp\{ M_{a,r,n}(t) \} dt = \frac{\sqrt{2\alpha}}{L} ( K(\nu) + o(1)),
$$
as $a \to 0$. 
The result now follows from \eqref{gacint}.
\end{proof}

\end{document}
